%-----------------------------------------------------------------------------------------------------------------------------------------------%
%	The MIT License (MIT)
%
%	Copyright (c) 2021 Jitin Nair
%
%	Permission is hereby granted, free of charge, to any person obtaining a copy
%	of this software and associated documentation files (the "Software"), to deal
%	in the Software without restriction, including without limitation the rights
%	to use, copy, modify, merge, publish, distribute, sublicense, and/or sell
%	copies of the Software, and to permit persons to whom the Software is
%	furnished to do so, subject to the following conditions:
%
%	THE SOFTWARE IS PROVIDED "AS IS", WITHOUT WARRANTY OF ANY KIND, EXPRESS OR
%	IMPLIED, INCLUDING BUT NOT LIMITED TO THE WARRANTIES OF MERCHANTABILITY,
%	FITNESS FOR A PARTICULAR PURPOSE AND NONINFRINGEMENT. IN NO EVENT SHALL THE
%	AUTHORS OR COPYRIGHT HOLDERS BE LIABLE FOR ANY CLAIM, DAMAGES OR OTHER
%	LIABILITY, WHETHER IN AN ACTION OF CONTRACT, TORT OR OTHERWISE, ARISING FROM,
%	OUT OF OR IN CONNECTION WITH THE SOFTWARE OR THE USE OR OTHER DEALINGS IN
%	THE SOFTWARE.
%
%
%-----------------------------------------------------------------------------------------------------------------------------------------------%

%----------------------------------------------------------------------------------------
%	DOCUMENT DEFINITION
%----------------------------------------------------------------------------------------

% article class because we want to fully customize the page and not use a cv template
\documentclass[a4paper,11pt]{article}

%----------------------------------------------------------------------------------------
%	FONT
%----------------------------------------------------------------------------------------

% codificacao de entrada/saida: necessario para acentuacao em portugues
\usepackage[utf8]{inputenc}
\usepackage[T1]{fontenc}
\usepackage{lmodern}

%----------------------------------------------------------------------------------------
%	PACKAGES
%----------------------------------------------------------------------------------------
\usepackage{url}
\usepackage{parskip}

%other packages for formatting
\RequirePackage{color}
\RequirePackage{graphicx}
\usepackage[usenames,dvipsnames]{xcolor}
\usepackage[scale=0.9]{geometry}

%tabularx environment
\usepackage{tabularx}

%for lists within experience section
\usepackage{enumitem}

% centered version of 'X' col. type
\newcolumntype{C}{>{\centering\arraybackslash}X}

%to prevent spillover of tabular into next pages
\usepackage{supertabular}
\usepackage{tabularx}
\newlength{\fullcollw}
\setlength{\fullcollw}{0.47\textwidth}

%custom \section
\usepackage{titlesec}
\usepackage{multicol}
\usepackage{multirow}

%CV Sections inspired by:
%http://stefano.italians.nl/archives/26
\titleformat{\section}{\Large\scshape\raggedright}{}{0em}{}[\titlerule]
\titlespacing{\section}{0pt}{10pt}{10pt}

%Setup hyperref package, and colours for links
\usepackage[unicode, draft=false]{hyperref}
\definecolor{linkcolour}{rgb}{0,0.2,0.6}
\hypersetup{colorlinks,breaklinks,urlcolor=linkcolour,linkcolor=linkcolour}

%for social icons
\usepackage{fontawesome5}

%debug page outer frames
%\usepackage{showframe}

% mes atual em portugues, sem depender do babel
\newcommand{\mesPT}{\ifcase\month\or janeiro\or fevereiro\or março\or abril\or maio\or junho\or julho\or agosto\or setembro\or outubro\or novembro\or dezembro\fi}


% job listing environments
\newenvironment{jobshort}[2]
    {
    \begin{tabularx}{\linewidth}{@{}l X r@{}}
    \textbf{#1} & \hfill &  #2 \\[3.75pt]
    \end{tabularx}
    }
    {
    }

\newenvironment{joblong}[2]
    {
    \begin{tabularx}{\linewidth}{@{}l X r@{}}
    \textbf{#1} & \hfill &  #2 \\[3.75pt]
    \end{tabularx}
    \begin{minipage}[t]{\linewidth}
    \begin{itemize}[nosep,after=\strut, leftmargin=1em, itemsep=3pt,label=--]
    }
    {
    \end{itemize}
    \end{minipage}
    }



%----------------------------------------------------------------------------------------
%	BEGIN DOCUMENT
%----------------------------------------------------------------------------------------
\begin{document}

% non-numbered pages
\pagestyle{empty}

%----------------------------------------------------------------------------------------
%	TITLE
%----------------------------------------------------------------------------------------

\begin{tabularx}{\linewidth}{@{} C @{}}
\Huge{Gustavo Fernandes Morais} \\[4pt]
\normalsize{Desenvolvedor Full Stack \ $\cdot$ \ Itabira, MG, Brasil} \\[7.5pt]
\href{https://github.com/GustavoF3rnandes}{\raisebox{-0.05\height}\faGithub\ GustavoF3rnandes} \ $|$ \
\href{https://www.linkedin.com/in/gustavo-fernandes-morais-6737a9243/}{\raisebox{-0.05\height}\faLinkedin\ gustavo-fernandes-morais} \ $|$ \
\href{https://portfoliomorais.vercel.app/}{\raisebox{-0.05\height}\faGlobe \ portfoliomorais.vercel.app} \ $|$ \
\href{mailto:gustavof3rnandes@outlook.com}{\raisebox{-0.05\height}\faEnvelope \ gustavof3rnandes@outlook.com} \ $|$ \
\href{tel:+5538998303009}{\raisebox{-0.05\height}\faMobile \ +55 (38) 99830-3009} \\[3pt]
% Links para as duas versoes do curriculo publicadas pelo GitHub Pages
\small{\raisebox{-0.05\height}\faLanguage \
\href{https://gustavof3rnandes.github.io/Curriculo/cv.pdf}{Português} \ $|$ \
\href{https://gustavof3rnandes.github.io/Curriculo/cv-en.pdf}{English}} \\
\end{tabularx}

%----------------------------------------------------------------------------------------
% EXPERIENCE SECTIONS
%----------------------------------------------------------------------------------------

%Interests/ Keywords/ Summary
\section{Resumo}
Engenheiro de Computação formado pela UNIFEI, atuando no ciclo completo de desenvolvimento de software: do front-end em React, Angular e Razor ao back-end em .NET, Python, Node.js e PHP, passando por modelagem de dados relacional e NoSQL, integrações via APIs REST e WebSocket e entrega em produção. Experiência consolidada em modernização de sistemas legados para arquiteturas atuais e na construção de aplicações serverless e de IA generativa na AWS, com 3 certificações AWS e 2 Google Cloud. Busco aprofundar conhecimentos em arquitetura de sistemas distribuídos, DevOps e aplicações de Inteligência Artificial.

%Experience
\section{Experiência Profissional}

\begin{joblong}{Cloud Engineer \& Solutions Architect, :upd8}{mar 2026 -- atual}
\item Modernização de aplicações legadas .NET para .NET Core 8 e 10 em clientes corporativos: análise do código-fonte existente, reescrita de controllers, views Razor e camada de acesso a dados, containerização com Docker e publicação via Azure Pipelines.
\item Modernização de plataforma contábil em PHP (TYPO3 Flow 2.3 / Doctrine) para Laravel 11 com Eloquent e arquitetura DDD, reorganizando aproximadamente 900 mil linhas de código legado em 20 domínios e 381 models.
\item Migração de bancos SQL Server para PostgreSQL com AWS DMS e Schema Conversion, incluindo reconstrução manual de foreign keys, triggers e stored procedures não convertidos automaticamente.
\item Uso de AWS Transform e Kiro para acelerar a refatoração de código legado, com revisão técnica das saídas e correção de regressões de segurança, roteamento e compatibilidade com Linux.
\item Definição e provisionamento de infraestrutura AWS segundo o Well-Architected Framework, do desenho de arquitetura à entrega em ambiente produtivo.
\item Condução de reuniões técnicas semanais com clientes, elaboração de documentação técnica e diagramas de arquitetura, além de apoio à frente comercial na construção de propostas.
\end{joblong}

\begin{joblong}{Desenvolvedor (Estagiário), :upd8}{abr 2025 -- mar 2026}
\item Desenvolvimento de aplicações full stack sob demanda para clientes e para processos internos: front-ends em React 19 com TypeScript e Vite consumindo APIs REST em Python (FastAPI) e Node.js.
\item Construção de APIs REST e WebSocket sobre Amazon API Gateway, com funções Lambda em Python, persistência em DynamoDB e processamento assíncrono via SQS.
\item Concepção e implementação de soluções de IA com Amazon Bedrock e, em fase subsequente, AWS AgentCore integrado ao framework Strands Agents, incluindo pipelines de RAG e agentes com ferramentas customizadas.
\item Automação de infraestrutura como código com AWS CDK, Terraform e AWS SAM, e pipelines de CI/CD em GitHub Actions.
\item Participação em reuniões de elicitação de requisitos funcionais e não funcionais, definição de escopo e elaboração de diagramas de arquitetura.
\item Conclusão de programa intensivo de capacitação em cloud computing com obtenção de 5 certificações profissionais (3 AWS e 2 GCP).
\end{joblong}

\begin{joblong}{Estagiário de Desenvolvimento Web e TI, ON Energia Solar}{set 2023 -- fev 2024}
\item Desenvolvimento e manutenção do site institucional utilizando WordPress e PHP.
\item Coleta e tratamento de dados com Microsoft Power BI para geração de indicadores à equipe.
\item Criação e gestão de campanhas no Google Ads e apoio na implementação de soluções internas de TI.
\end{joblong}

\begin{joblong}{Desenvolvedor Front-end, Prefeitura Municipal de Itabira}{jul 2022 -- jun 2023}
\item Desenvolvimento do Portal de Turismo de Itabira com ASP.NET Razor, criando as interfaces da aplicação com foco em acesso multiplataforma.
\item Implementação das principais funções de negócio em C\# e das interações com banco de dados SQL Server.
\item Trabalho em equipe com práticas de web design e gerenciamento de projeto junto ao órgão público.
\end{joblong}

%Projects
\section{Projetos}

\begin{jobshort}{AI Solutions Advisor}{\textit{React 19, TypeScript, FastAPI, Bedrock}}
Aplicação full stack de chat com agente de IA. Front-end em React 19 com TypeScript e Vite, com interface de conversa e painel de debug das chamadas de ferramenta; back-end em Python com FastAPI e validação via Pydantic, expondo os endpoints \texttt{/api/chat} e \texttt{/health}. O agente Strands sobre Amazon Bedrock dispõe de ferramentas customizadas de cálculo, data/hora e consulta a catálogo. Testes com Vitest e Testing Library no front-end e pytest com Hypothesis no back-end.
\end{jobshort}

\begin{jobshort}{Modernização do sistema PDPoint}{\textit{ASP.NET Core 8, C\#, PostgreSQL, Docker}}
Migração de aplicação ASP.NET MVC em .NET Framework para ASP.NET Core 8: mais de 20 controllers e 180 views Razor reescritos, troca da camada de dados para Npgsql (PostgreSQL), logging estruturado com Serilog, mapeamento com AutoMapper, geração de PDF com PuppeteerSharp e containerização para deploy via Azure Pipelines.
\end{jobshort}

\begin{jobshort}{APIs de Portabilidade e Facade}{\textit{Node.js, TypeScript, Express, Prisma}}
Modernização de APIs legadas de uma operadora de telecomunicações. A API Facade foi reescrita em Node.js 22 com TypeScript, Express e Prisma, organizada em Clean Architecture (entities, use cases, handlers), com integrações resilientes usando retry e circuit breaker, testes em Vitest, secrets e storage no Google Cloud e deploy automatizado. A API de Portabilidade foi modernizada em .NET com jobs em background, webhooks e execução em Kubernetes.
\end{jobshort}

\begin{jobshort}{Plataforma multi-tenant de agentes conversacionais}{\textit{AWS CDK, Python, Bedrock}}
Chatbot integrado ao WhatsApp com isolamento completo por cliente. Webhooks chegam via API Gateway e acionam uma Lambda de ACK que consolida fragmentos de mensagem em buffer no DynamoDB e enfileira no SQS com delay; a Lambda de invocação recupera dinamicamente as instruções do cliente e responde pelo Bedrock com Knowledge Base própria. EventBridge mantém a base vetorial sincronizada e cada transação registra tracing e consumo de tokens para rateio de custo. Infraestrutura inteiramente em AWS CDK (Python), com deploy por GitHub Actions.
\end{jobshort}

\begin{jobshort}{Validação e extração inteligente de documentos}{\textit{Python, Lambda, Textract, Terraform}}
Pipeline serverless que valida autenticidade e extrai campos estruturados de 12 grupos de documentos brasileiros, combinando Rekognition, Textract e Bedrock. Arquitetura modular com validators e extractors dedicados por tipo de documento aplicando Template Method, API REST documentada em OpenAPI, Lambda Layers para dependências e infraestrutura provisionada com Terraform.
\end{jobshort}

\begin{jobshort}{Projetos pessoais}{\href{https://github.com/GustavoF3rnandes}{\textit{github.com/GustavoF3rnandes}}}
Notes-React (React + TypeScript), Food-ReactNative (React Native + TypeScript), Simulador de Bancada de Incêndios (C\#), tratamento de dados de planilhas em Python e portfólio pessoal publicado na Vercel.
\end{jobshort}

%----------------------------------------------------------------------------------------
%	EDUCATION
%----------------------------------------------------------------------------------------
\section{Formação Acadêmica}
\begin{tabularx}{\linewidth}{@{}l X@{}}
2021 -- 2026 & Bacharelado em Engenharia de Computação na \textbf{Universidade Federal de Itajubá (UNIFEI)} \\
\end{tabularx}

%----------------------------------------------------------------------------------------
%	CERTIFICATIONS
%----------------------------------------------------------------------------------------
\section{Certificações}
\begin{tabularx}{\linewidth}{@{}X X@{}}
AWS Certified Developer -- Associate & Google Professional Cloud Architect \\
AWS Certified Advanced Networking -- Specialty & Google Associate Cloud Engineer \\
AWS Certified Cloud Practitioner & \\
\end{tabularx}

%----------------------------------------------------------------------------------------
%	SKILLS
%----------------------------------------------------------------------------------------
\section{Competências Técnicas}
\begin{tabularx}{\linewidth}{@{}l X@{}}
Front-end & \normalsize{React, React Native, TypeScript, JavaScript, Angular, Razor / ASP.NET MVC, HTML5, CSS3, Bootstrap, Vite} \\
Back-end & \normalsize{C\# e .NET (Framework, 8 e 10), Python (FastAPI), Node.js (Express), PHP (Laravel, WordPress)} \\
Dados & \normalsize{SQL Server, PostgreSQL, DynamoDB, Prisma, Eloquent, Entity Framework, AWS DMS} \\
IA aplicada & \normalsize{Amazon Bedrock, Strands Agents, AWS AgentCore, RAG e Knowledge Bases, Textract, Rekognition} \\
Cloud \& DevOps & \normalsize{AWS, Google Cloud, Docker, Kubernetes, Terraform, AWS CDK, AWS SAM, GitHub Actions, Azure DevOps} \\
Práticas & \normalsize{APIs REST e WebSocket, OpenAPI, Clean Architecture, DDD, testes automatizados (pytest, Vitest, Testing Library), Git, levantamento de requisitos} \\
Idiomas & \normalsize{Português (nativo), Inglês (profissional), Espanhol (básico)} \\
\end{tabularx}

\vfill
\center{\footnotesize Última atualização: \mesPT{} de \the\year}

\end{document}
